\section{Задание}

\textbf{Цель работы:} освоение принципов проектирования многотактных цифровых устройств на языке Verilog HDL, реализация конечных автоматов с раздельным описанием логики переходов и регистров состояния, интеграция арифметических модулей, а также проведение полного цикла проектирования для ПЛИС на базе отладочной платы Saylinx.

В процессе выполнения лабораторной работы необходимо:
\begin{itemize}
    \item Разработать и описать на языке Verilog HDL многотактные модули арифметических операций с сигналами управления \texttt{start} и \texttt{done}.
    \item Реализовать управление вычислениями в виде конечных автоматов с явным разделением на комбинационную часть логики следующего состояния и последовательностную часть памяти состояния.
    \item Разработать многотактную схему вычисления целевой функции, объединяющую созданные арифметические модули в соответствии с заданными ограничениями и разрядностью данных (16-битные входные аргументы).
    \item Определить область допустимых значений реализуемой функции.
    \item Построить граф переходов для FSM каждого разработанного модуля.
    \item Разработать тестовое окружение для проверки работоспособности схемы не менее чем на 10 различных тестовых векторах.
    \item Провести функциональное моделирование с учетом частоты тактирования 100 МГц и определить время выполнения вычислений.
    \item Выполнить синтез и размещение схемы для ПЛИС платы Saylinx, проанализировать аппаратные затраты.
    \item Подготовить полный отчет по результатам выполнения работы.
\end{itemize}

Вариант работы определяется последней цифрой номера ИСУ студента. В моём случае вариант работы равен:
\[
    \texttt{get\_variant}(466449) = 9.
\]

Этому варианту соответствует функция:
\[
    y=\sqrt[3]{a}+\sqrt{b}
\]

Она обладает следующими ограничениями:
\begin{itemize}
    \item количество сумматоров --- 1;
    \item количество умножителей --- 2.
\end{itemize}